\chapter{Introduction}
In the last years, the increasing availability of connected devices has progressively shifted the data generation process from centralized infrastructures towards large populations of edge devices, including smartphones, wearables, vehicles and sensors.
According to~\cite{cisco2023annual}, the number of networked devices and connection was expected to reach 29.3 billion by 2023, corresponding to more than three times the global population, and approximately 3.6 devices per person. Each of these devices, represents a valuable data generation source, collecting infromation about daily interaction between individuals, devices and digital services.

In this context, Machine Learning (ML) emerges as an incredibly powerful technology to automatically extract patterns from enormous amount of data. In general, access to more data allows to capture richer statistical regularities, improve generalization, and adapt to a wider range of users, environments and tasks~\cite{banko-brill-2001-scaling, shalev2014understanding, hestness2017deeplearningscalingpredictable, kaplan2020scalinglawsneurallanguage}. 

However, having access to more data does not necessarly imply that these data can be freely collected or used. In many domains, data are subject to legal, ethical, technical, and social constraints~\cite{eu_gdpr_2016, hiroshima_process_guiding_principles_2023, eu_ai_act_2024, Buick_copyright, eu_ehds_2025, california_ccpa_cpra}. Medical records are protected because they describe a person’s health status. Financial data may reveal income, consumption patterns, debt, or economic vulnerability. Text, images, and user-generated content may be protected by copyright or by terms of service. Location traces may disclose where a person lives, works, studies, or receives medical care. Even apparently harmless digital traces can become sensitive when aggregated over time or combined with other sources of information \cite{netflix_attack}. Data are therefore not only an input to optimization procedures but are also representations of individuals, their habits, preferences, behaviors, and social relations. 

This observation raises a broader question about control. If data generated by personal devices can reveal intimate aspects of people’s lives, who has access to such data, and under which conditions? In practice, large private companies often occupy a privileged position in the digital ecosystem. They design the devices, operating systems, applications, cloud services, and platforms through which data are generated, stored, and processed. As a consequence, they may acquire detailed knowledge about users’ behavior, preferences, movements, purchases, communications, and interactions. This knowledge can be used in many useful ways to improve people’s daily living, but at the same time, it can also be used for profiling, targeting, ranking, persuasion, facial recognition, or other forms of automated decision-making that are not always transparent to the individuals concerned and are beneficial only for few peoples \cite{pasquale_black_box_society_2015, veliz_privacy_is_power_2021}.

The question is therefore not only whether data are useful, but whether the entities that collect and process them can be trusted. Recent enforcement actions involving major technology companies suggest that privacy risks often arise not only from external attacks or accidental data breaches, but also from ordinary platform practices: collecting data under opaque conditions, reusing it for purposes different from those initially presented to users.

The Cambridge Analytica scandal showed how personal information collected through a digital platform could be repurposed for political profiling and targeted persuasion, raising public concern about consent, transparency, and the secondary use of personal data \cite{ico_democracy_disrupted_2018, ico_data_analytics_political_campaigns_2018}. This episode is part of a broader pattern of documented privacy controversies involving major technology companies \cite{cnil_google_gdpr_2019, ftc_youtube_coppa_2019, ftc_twitter_security_data_ads_2022, ftc_amazon_alexa_coppa_2023, texas_ag_google_privacy_2025}. These cases suggest that privacy risks do not arise only from external breaches, but also from the ordinary ways in which platform operators collect, retain, combine, and repurpose personal data.

Similar concerns arise in machine learning, where the demand for large and diverse training datasets creates strong incentives to collect and reuse data at scale. Controversies involving biometric systems trained on scraped or improperly obtained facial images \citeweb{noyb_clearview_criminal_complaint_2025}, copyrighted material used for model training \citeweb{reisner_ai_pirated_books_problem_2025}, and wearable devices capturing data from both users and bystanders \citeweb{svd_meta_ai_glasses_privacy_2026} illustrate how the search for training data can blur the boundary between legitimate data use and unauthorized appropriation.


These concerns highlight two obstacles to the conventional centralized learning paradigm. The first is practical. Data are increasingly generated across large populations of devices, so collecting them into a single repository may require substantial communication and storage costs, making centralized collection inefficient and difficult to maintain at scale. 

The second obstacle concerns privacy and trust. Even when centralization is technically feasible, transferring raw records to a server increases the exposure of sensitive information and requires users or institutions to trust the entity that collects, stores, and processes the data. The problem is therefore not only how to learn from more data, but how to learn from distributed data without requiring all raw records to be pooled under the control of a single actor.

Federated Learning (FL) has emerged as a response to this challenge \cite{konecny2015federatedoptimizationdistributedoptimizationdatacenter, pmlr-v54-mcmahan17a}. Instead of collecting raw data in a central repository, FL moves the training procedure closer to the data sources: clients keep their datasets locally, compute model updates on their own devices or institutional servers, and communicate only these updates to a central coordinator. In this way, FL addresses the distributed nature of modern data generation while reducing the need to entrust a central entity with direct access to raw personal records.



